%--------------------
% Knowledge Graphs (VU 192.116, 2026S) -- Portfolio Report
% A Temporal Football Knowledge Graph for Tactical Pattern Discovery
% and Team Style Comparison
% Tim Gress (12412672), Jan Toelken (12432831) -- 6 ECTS
%
% Structure follows the official "Portfolio Example Structure" 1:1.
% -> final PDF filename must end in "-structured.pdf"
% -> ZIP folders: "2 - construction", "3 - ML", "4 - logic", "5 - reflection"
% TODO markers mark what to fill in / produce before submission.
%--------------------

\documentclass[11pt,a4paper]{article}
\usepackage[utf8]{inputenc}
\usepackage[T1]{fontenc}
\usepackage{mathptmx}                 % Times font

\usepackage[pdftex]{graphicx}
\usepackage[pdftex,linkcolor=black,pdfborder={0 0 0}]{hyperref}
\usepackage{booktabs}                 % clean tables
\usepackage{enumitem}
\usepackage{verbatim}

\frenchspacing
\linespread{1.15}
\usepackage[a4paper,margin=2.5cm]{geometry}
\usepackage[all]{nowidow}
\usepackage[protrusion=true,expansion=true]{microtype}

\hypersetup{
  pdftitle  = {A Temporal Football Knowledge Graph},
  pdfauthor = {Tim Gress, Jan Toelken}
}

\title{A Temporal Football Knowledge Graph for\\
       Tactical Pattern Discovery and Team Style Comparison}
\author{Tim Gre\ss{} (12412672) \and Jan T\"olken (12432831)}
\date{Knowledge Graphs (VU 192.116, 2026S) \quad--\quad 6 ECTS
      \quad--\quad [TODO date]}

\begin{document}
\maketitle

% TODO: insert the filled cover pages (pro forma) BEFORE this page in the final
% PDF. Brief per-LO formulations are in report/cover-pages-notes.md.

\begin{abstract}
\noindent
We build a temporal knowledge graph from StatsBomb open football event data
(Premier League 2015/16; 380 matches, 1.31 million events), formalise three
tactical patterns -- pressing regains, fast transitions and wide build-up -- as
graph queries and recursive Datalog rules, and use the materialised pattern
instances to compare the playing styles of all 20 teams through an interactive
analytics dashboard. A knowledge-graph-embedding experiment (TransE, ComplEx)
complements the symbolic pipeline; a multi-seed evaluation reproduces the
expected model contrast on 1-to-N relations but shows that, at this graph size,
the learned space does not reliably recover the rule-derived style structure.
All code and data references are in the accompanying ZIP file.
\end{abstract}

%--------------------
\section{Scenario}
%--------------------

\subsection{Application domain: football analytics}

Professional football produces fine-grained event data: every pass, pressure,
carry and shot of a match is recorded with player, team, location and timestamp.
Clubs, media and analytics providers use this data to answer \emph{tactical}
questions -- does a team press high, counter-attack, build up through the wings?
-- questions that are inherently about \emph{sequences of connected events}
rather than isolated statistics. This makes the domain a natural fit for a
knowledge graph: events are related temporally (what happened next), spatially
(where on the pitch) and organisationally (who did it, for which team), and
tactical concepts are patterns over exactly these relations (LO9). Comparable
industrial systems exist at football-analytics companies and clubs' data
departments; our project realises the same idea end to end at portfolio scale
(LO9).

\subsection{The service provided through the KG}

The service is \emph{tactical pattern discovery and team-style comparison}
(LO11): (i) a class of graph queries/rules that finds instances of three
tactical patterns in the event stream and materialises them back into the graph
as derived knowledge; and (ii) an interactive dashboard that aggregates these
instances into per-team style profiles -- pressing intensity, directness, wide
play, possession behaviour -- and lets a user explore the league, a single team,
or a head-to-head comparison, down to single reconstructed goal sequences. The
outcome of the service is evaluated in Section~\ref{sec:service-outcome}.

%--------------------
\section{KG Construction}
%--------------------

\subsection{Datasets}

We use StatsBomb Open Data,\footnote{\url{https://github.com/statsbomb/open-data}
-- used under the StatsBomb non-commercial licence; StatsBomb is credited as the
data source.} the largest publicly available collection of professional football
event data. Each match is a JSON file of about 3\,500 events; every event carries
common fields (\texttt{id}, \texttt{index}, \texttt{period}, \texttt{timestamp},
\texttt{type}, \texttt{possession}, \texttt{team}, \texttt{player},
\texttt{location}) plus type-specific detail objects. Across the season the data
contains \textbf{33 distinct event types}; the most frequent are Pass (28.1\%),
Ball Receipt (25.9\%), Carry (21.1\%) and Pressure (8.8\%).

From the roughly 12\,GB collection we selected the \textbf{Premier League
2015/16 season} (380 matches). Profiling showed it is one of the only complete,
balanced league seasons in the open data (every team plays 38 matches), whereas
most other competitions are single-club-centric; a balanced league is essential
for a fair \emph{comparison} of teams. The season is also the one of Leicester
City's counter-attacking title win, which -- as shown later -- our patterns
rediscover from the raw event stream.

\subsection{Technologies and architecture}

We store the KG as a property graph in \textbf{Neo4j~5.26} (Docker, APOC
plugin); the data-model rationale is discussed in
Section~\ref{sec:datamodel}. A thin Python layer surrounds it: ingestion (JSON
to Neo4j, batched, one transaction per match, resumable and idempotent); the
pattern layer (Cypher path queries and recursive Datalog rules); the service
layer (Streamlit dashboard); and the embedding layer (PyKEEN, in a separate
environment). The division of labour follows the course principle: the KG
handles \emph{structure and derivation} -- temporal order, containment, pattern
instances as first-class derived nodes -- while application code handles
presentation, statistics and ML training (LO5). Everything a later query must
traverse is an edge; everything only rendered is computed outside the graph.

% TODO: Figure 1 -- architecture/pipeline diagram
% (StatsBomb JSON -> ingest -> Neo4j KG -> {patterns/Datalog, dashboard,
%  triple export -> PyKEEN}). Reference it here.

\subsection{Constructing the KG: schema design and mapping}
\label{sec:construction}

This is one of the two focus areas of the project (LO7).

\paragraph{Design goals.} The schema is driven by what the patterns need:
temporal order (patterns are sequences), possession grouping (patterns live
inside possession chains), spatial context (patterns reference pitch areas) and
participation (per-team aggregation).

\paragraph{Schema.} Every event becomes an \texttt{:Event} node; subtype labels
(\texttt{:Pass}, \texttt{:Shot}, \texttt{:Pressure}, \dots) carry type-specific
properties. Temporal order is materialised as an explicit
\texttt{(:Event)-[:NEXT]->(:Event)} chain per match -- the single most important
design decision, because it turns ``within $k$ steps afterwards'' into a
variable-length path. Possessions are reified as \texttt{:Possession} nodes
(events linked via \texttt{[:PART\_OF]}); a static $3\times3$ \texttt{:Zone}
grid gives spatial context via \texttt{[:IN\_ZONE]}; \texttt{[:BY\_TEAM]},
\texttt{[:BY\_PLAYER]} and \texttt{[:RECEIVED\_BY]} capture participation. Source
coordinates are already oriented per acting team (x $\to$ 120 = opponent goal),
so ``own half'' is uniformly $x < 60$.

\paragraph{Three concrete mapping examples.}
\begin{enumerate}[itemsep=2pt]
  \item \emph{Event node with subtype.} A JSON pass maps to one
    \texttt{(:Event:Pass)} node with common properties plus pass-specific ones
    (\texttt{length}, \texttt{outcome}, \texttt{end\_x/y}); StatsBomb omits
    \texttt{outcome} for completed passes, so the mapping defaults it to
    \texttt{"Complete"} -- a small documented cleaning decision.
  \item \emph{NEXT edge.} Consecutive \texttt{index} values within a match map
    to a \texttt{[:NEXT]} edge; verification asserts the chain is a single path
    ($n-1$ edges for $n$ events, one start per match).
  \item \emph{Possession node (reification).} The StatsBomb possession counter
    (an integer on each event) is reified into a \texttt{:Possession} node keyed
    on \texttt{(match\_id, possession)}, turning an implicit grouping into a
    first-class object that patterns anchor on.
\end{enumerate}

\paragraph{Identity and record linkage.} StatsBomb IDs are globally consistent
across files, so no record linkage was required for this single-source KG -- the
classically hard part of KG creation does not arise here, which we state rather
than manufacture (LO7). Linking players/teams to Wikidata is noted as future
work.

\paragraph{Verification.} Ingestion is time-boxed; correctness is not. A
verification script reconciles the KG against the source (per-label node counts,
per-match event counts equal to the JSON length, NEXT-chain integrity, every
event in exactly one possession, every team with 38 matches). Final contents:
\textbf{1\,313\,773 events} (368k passes, 277k carries, 115k pressures, 9\,908
shots, \dots), \textbf{71\,884 possessions}, 644 players, 380 matches, 20 teams,
9 zones.

%--------------------
\section{ML-based Representation: KG Embeddings}
%--------------------

\subsection{The representation and its training}

We train knowledge-graph embeddings with PyKEEN on triples exported from the KG
(LO1). Embeddings need a discrete entity--relation structure; the raw event
stream with continuous coordinates does not embed meaningfully (see
Section~\ref{sec:datamodel}), so we export the aggregated player/team layer:
\textbf{9\,300 triples over 596 entities and 4 relations} --
\texttt{plays\_for} (561), \texttt{plays\_position} (2\,179), \texttt{passes\_to}
(5\,774; completed passes between player pairs with at least 10 occurrences) and
\texttt{exhibits\_pattern} (786; players in at least 10 pattern-instance
anchors, feeding derived knowledge from Section~\ref{sec:logic} into the
embedding). We train TransE (translation-based) and ComplEx (semantic matching),
$d = 64$, 150 epochs, 80/10/10 split, filtered evaluation. Because the graph is
small, we repeat the whole experiment over \textbf{eight random seeds} (both the
split and the initialisation) and report mean $\pm$ standard deviation rather
than a single run:

\begin{table}[ht]
\centering
\begin{tabular}{lcccc}
\toprule
Model & MRR & Hits@1 & Hits@3 & Hits@10 \\
\midrule
TransE  & $0.266 \pm 0.010$ & $0.037 \pm 0.009$ & $0.404 \pm 0.021$ & $0.722 \pm 0.015$ \\
ComplEx & $0.321 \pm 0.025$ & $0.176 \pm 0.029$ & $0.395 \pm 0.029$ & $0.608 \pm 0.023$ \\
\bottomrule
\end{tabular}
\caption{Link prediction on held-out test triples, mean $\pm$ standard
deviation over eight seeds.}
\end{table}

The contrast is textbook and robust across every seed (LO1): TransE's near-zero
Hits@1 against ComplEx's 0.176 reflects the known weakness of translational
models on 1-to-N relations -- a team has many players, a position many holders,
so $h + r \approx t$ cannot hold for all tails at once, while ComplEx's bilinear
scoring can. The models fail in different directions, which the two ends of the
ranking make visible: TransE is the \emph{better} model at Hits@10 ($0.722$ vs.\
$0.608$) but five times worse at Hits@1 -- its translation places the correct
entity in the right region of the space without being able to single it out,
exactly the behaviour the 1-to-N geometry predicts.

% TODO: 5 concrete example representations (e.g. nearest neighbours of Leicester
% City, Kante, the P1 pattern node, a goalkeeper, a relation) from
% generated/embeddings/team_neighbours_complex.txt.
% TODO: one true-positive and one false-positive predicted link (small script).

\subsection{Use for evolving the KG: completion}

The held-out evaluation \emph{is} KG completion (LO8): hiding true edges and
recovering them simulates an incomplete graph. Hits@10 $= 0.72$ (TransE, stable
across seeds) means a missing link is recovered among the top 10 candidates 72\%
of the time -- usable
as a suggestion mechanism (e.g.\ likely pass partnerships for a new player) with
a human or threshold in the loop, not an automatic writer. We use it in the
completion sense; cleaning/correcting with embeddings is out of scope.

\subsection{Context, limitations, scalability}

Four honest limitations (LO6): (i) \emph{volume} -- 9.3k triples is small;
embeddings shine at millions, and our export trades volume for semantic density.
(ii) \emph{cold start} -- entities without training triples (new players) get no
meaningful vector. (iii) \emph{scaling} -- training was CPU-only in minutes and
scales on GPU; the binding constraint is not compute but \emph{what to export},
since literals, timestamps and coordinates do not embed in this model family --
exactly the trade-off in Section~\ref{sec:datamodel}. (iv) \emph{variance} --
and this one changed a conclusion of ours. Single-seed results looked
convincing until we repeated the experiment: the link-prediction metrics are
stable (standard deviations of $0.01$--$0.03$), but any statistic computed over
only the 20 team entities swings widely between seeds. The downstream
cross-check in Section~\ref{sec:connections} varied from 4 to 11 out of 20
across the eight runs -- a single run would have supported the opposite
conclusion of the one we now draw. At this graph size, multi-seed reporting is
not optional rigour but the difference between a true and a false statement.

%--------------------
\section{Logic-based Representation: Tactical Patterns as Rules}
\label{sec:logic}
%--------------------

\subsection{The representation: queries and recursive rules}

Tactical patterns are formalised twice over the same KG (LO2): as Cypher path
queries (production path) and as recursive Datalog rules (Soufflé; executable
equivalently as recursive SQL CTEs in DuckDB on exported facts). Five examples,
one recursive and one creating objects:

\begin{enumerate}[itemsep=3pt]
  \item \emph{P1 -- Pressing Regain (Cypher, recursive).} A pressure in the
    opponent half followed within 5\,s by a regain of the pressing team, as a
    variable-length path over the temporal chain:
\begin{verbatim}
MATCH (p:Pressure)-[:BY_TEAM]->(t:Team) WHERE p.x >= 60
MATCH (p)-[:NEXT*1..10]->(r:Event)
WHERE r.timestamp_s - p.timestamp_s <= 5 AND (r)-[:BY_TEAM]->(t)
  AND (r:BallRecovery OR r:Interception OR r:Duel ...)
\end{verbatim}
    \texttt{NEXT*1..10} is full recursion over the event chain, bounded by the
    time window.
  \item \emph{P1 in Datalog (the recursive core).} The transitive closure of
    \texttt{next}, bounded by the 5-second window:
\begin{verbatim}
reach(P,E,1)   :- pressure(P,_,_,TS,PD), next(P,E),
                  event(E,_,_,PD,TE,...), TE-TS <= 5.
reach(P,E,D+1) :- reach(P,M,D), D < 10, next(M,E),
                  pressure(P,_,_,TS,PD),
                  event(E,_,_,PD,TE,...), TE-TS <= 5.
p1(P) :- pressure(P,_,T,_,_), reach(P,R,_), wins_ball(R,T).
\end{verbatim}
  \item \emph{P2 -- Fast Transition.} An open-play possession whose first
    located event lies in the own half and which produces a shot within 15\,s
    and at least 30 units of forward progress (non-recursive join over
    possession containment).
  \item \emph{P3 -- Wide Build-Up.} A Regular-Play possession with at least 3
    own events in the wide middle-third zones, later entering the wide final
    third -- spatial reasoning via the \texttt{:Zone} abstraction.
  \item \emph{Materialisation as object creation.} Every match writes a new
    \texttt{(:PatternInstance)} node with \texttt{[:MATCHES]} provenance edges
    to its events, \texttt{[:FOUND\_IN]} to the match and
    \texttt{[:EXHIBITED\_BY]} to the team -- new nodes derived by rules, the
    object-creation aspect of LO2 (idempotent via \texttt{MERGE} on a
    pattern/match/anchor key).
\end{enumerate}

Over the full season this yields \textbf{9\,121 P1, 979 P2 and 6\,998 P3
instances}. Two validations support the definitions. (i) \emph{Vendor signals}:
P2 instances are about 17$\times$ enriched for StatsBomb's own
\texttt{From Counter} label (30.5\% vs.\ a 1.8\% base rate), capturing 69.7\% of
counter-labelled possessions that contain a shot; P1 overlaps the
\texttt{counterpress} flag only weakly (39.6\% vs.\ 36.9\% baseline), as
expected -- the flag marks pressing after \emph{own} loss while P1 requires the
press to \emph{succeed} -- which we report honestly rather than tune towards.
(ii) \emph{Cross-formalism}: the Datalog rules reproduce the Cypher counts
\emph{exactly} (9\,121 = 9\,121 and 979 = 979, per match, zero deviations),
evidence that the formal definitions, not implementation accidents, determine
the result (LO2).

\subsection{Use for evolving the KG: derived knowledge}

The rules \emph{add} to the graph (LO8): 17\,098 \texttt{PatternInstance} nodes
and their provenance edges are new knowledge derived from, and stored alongside,
the base data. Downstream consumers -- the style metrics, the dashboard, even
the \texttt{exhibits\_pattern} triples of the embedding export -- read the
derived layer without re-running the rules. Schema evolution is cheap in the
property-graph model: new event types already enter as plain \texttt{:Event}
nodes (keeping the NEXT chain intact) and can be upgraded with subtype labels
later without re-ingestion; adding a season is a re-run of the idempotent loader
(LO8).

\subsection{Context, limitations, scalability}

Scalable reasoning was not free (LO6). During ingestion, an event lookup through
a \emph{subtype} label bypassed the unique index defined on
\texttt{:Event(event\_id)} -- Neo4j indexes are label-specific -- causing full
label scans and a quadratic slowdown (stalling around match 80 of 380); the
corrected lookup restored about 5\,000 events/s. The pattern queries then
initially label-scanned the whole season per match: \texttt{PROFILE} showed the
composite \texttt{:Event(match\_id, idx)} index never applied, because composite
indexes require predicates on all their properties. Single-property
\texttt{match\_id} indexes plus one \texttt{USING INDEX} hint made P1
4$\times$, P2 4$\times$ and P3 13$\times$ faster at identical results. As a
comparison point, the same patterns as recursive SQL CTEs in DuckDB run in about
1\,s over the whole season on columnar exports -- fast, but the rules read as
machinery (explicit fixpoint construction) rather than the tactical statement
the Cypher path expresses; discussed further in
Section~\ref{sec:connections}.

%--------------------
\section{Reflection}
%--------------------

\subsection{Outcome of the service}
\label{sec:service-outcome}

The service was delivered (LO11): from the materialised pattern instances we
compute five per-team style metrics -- each normalised \emph{per opportunity}
(e.g.\ pressing per opponent possession, directness per own-half open-play
possession) so that possession imbalances do not distort comparison -- and
provide them through an interactive Streamlit dashboard backed by the KG:
league overview (final table with style profile, style map, automatically
clustered ``style families''), team profiles (auto-generated plain-language
verdicts, rank metrics, pitch maps on a league-wide shared colour scale) and
head-to-head comparison; drill-downs reconstruct single goal moves by walking a
possession's event chain in the graph.

The headline result is the \emph{face validity} of the discovered styles.
Leicester City -- the counter-attacking champions of exactly this season -- rank
first in directness by a wide margin (0.052 vs.\ 0.038 for the next team); the
renowned pressing sides (Tottenham, Liverpool, Manchester City) top pressing
success (23.0\% / 21.2\% / 22.7\% of high presses regain the ball within 5\,s
vs.\ a 19.6\% league average); van Gaal's Manchester United tops wide build-up.
None of this football knowledge was given to the system -- it emerges from rules
over the graph.

% TODO: 1-2 dashboard screenshots + the style-metric table as figures.

Limitations: patterns are threshold-based (a sensitivity analysis is future
work); one season, one league (the pipeline generalises -- a second 2015/16
league would enable cross-league comparison); and linkage to external KGs
remains the natural next construction step.

\subsection{Data model: comparison and reflection}
\label{sec:datamodel}

We compared four data models before committing (LO4). \emph{Relational tables}
are the source's natural shape, but sequence patterns become self-join chains
and variable-length sequences need recursive CTEs -- our DuckDB re-implementation
shows both feasibility and awkwardness. \emph{RDF} would explode each
$\sim$10-attribute event into many triples and needs reification (or RDF-star)
for time on statements. \emph{Temporal KG models} with interval-annotated edges
suit slowly-changing facts (transfers, contracts), not a stream of instantaneous
events. We chose a \emph{temporal property graph}: an event is one node with
properties; time is both a property (for $\Delta t$ arithmetic) and structure
(the NEXT chain); sequences are paths; patterns map 1:1 to Cypher and to Datalog
facts. The embedding experiment adds the ML perspective: to become vectors, the
property graph had to be flattened to id-triples, discarding literals and time
-- each representation keeps what it can use (LO4).

\subsection{Connections between KGs, ML and AI}
\label{sec:connections}

Our project instantiates the course's ``melting pot'' claim on one concrete KG
(LO12). The symbolic side (recursive rules) and the ML side (embeddings) not
only coexist -- they feed each other: pattern instances derived by logic become
training triples for the embeddings (\texttt{exhibits\_pattern}), and the
learned space is then evaluated against the symbolic style clusters: for each
team we ask whether its nearest neighbour in embedding space falls in the same
Ward cluster that the rule-derived style metrics produce. The chance level for
this test is \emph{not} $1/4$: the clusters have unequal sizes, so the expected
number of agreements under random neighbours is
$\sum_t (|c_t| - 1)/(N-1) = 6.0$ of 20.

Measured over eight seeds, ComplEx reaches $7.8 \pm 2.6$ agreements and TransE
$4.9 \pm 1.4$. Neither is distinguishable from chance (one-sample $t$-test
against $6.0$: ComplEx $t = 1.90$, $p = 0.10$; TransE $t = -2.35$, $p = 0.05$,
if anything \emph{below} chance). We therefore cannot claim that the
sub-symbolic representation recovers the symbolic style structure -- a claim a
single lucky seed ($9/20$) would have supported. What the experiment does show
is more modest and more useful: the two representations encode \emph{different}
things about the same graph. The embeddings are trained on passing topology,
positions and pattern participation, and they predict links well
(Hits@10 $\approx 0.72$); playing style as we define it is a normalised,
per-opportunity aggregate that this topology apparently does not determine.

The comparison also shows a clean division of strengths (LO12): logic gave
exact, explainable, provably equivalent results (Cypher $\equiv$ Datalog to the
instance); embeddings gave generalisation to unseen edges at the cost of
explainability, and -- as above -- of stability at this scale; the graph
database contributed the systems dimension (indexes, query plans) that made both
feasible. A larger export (multiple seasons, player-level pattern relations) is
the obvious way to give the sub-symbolic side a fair chance, and rule learning
-- inducing the pattern definitions from data instead of writing them -- would
close the remaining loop between the two sides.

%--------------------
\section*{Attribution}
%--------------------
Data: StatsBomb Open Data, used under their non-commercial licence; StatsBomb is
credited as the source of all event data. % TODO: StatsBomb logo in data figures.
% TODO: reuse declaration on the cover pages (declare AI-assisted development).

%--------------------
\appendix
\section{Reproducibility}
%--------------------
The ZIP file mirrors this report's structure (\texttt{2 - construction},
\texttt{3 - ML}, \texttt{4 - logic}, \texttt{5 - reflection}), each folder with
a \texttt{readme.md}. The full dataset exceeds the upload limit; the report links
the public StatsBomb repository and the ZIP contains a two-match subset for
partial reproducibility. One command per stage: \texttt{docker compose up -d};
\texttt{python -m src.ingest.load} and \texttt{verify};
\texttt{python -m src.patterns.run} and \texttt{validate};
\texttt{python -m src.datalog.export} and \texttt{run\_duckdb};
\texttt{python -m src.dashboard.metrics} and
\texttt{streamlit run src/dashboard/app.py};
\texttt{python -m src.embeddings.export\_triples} and \texttt{train}.

\end{document}
